\documentclass[11pt, a4paper]{report}

% ========== Common Packages ==========
\usepackage{fontspec}
\usepackage{kotex}
\usepackage{amsmath, amssymb, amsthm}
\usepackage{mathtools}
\usepackage{bm}
\usepackage{graphicx}
\usepackage{booktabs}
\usepackage{array}
\usepackage{tabularx}
\usepackage{longtable}
\usepackage{hyperref}
\usepackage{cleveref}
\usepackage{geometry}
\usepackage{fancyhdr}
\usepackage{titlesec}
\usepackage{enumitem}
\usepackage{xcolor}
\usepackage{tcolorbox}
\usepackage{algorithm}
\usepackage{algorithmic}
\usepackage{tikz}
\usetikzlibrary{arrows.meta, positioning, calc, decorations.pathreplacing, shapes.geometric, fit, patterns, angles, quotes, trees}
\usepackage{pgfplots}
\pgfplotsset{compat=1.18}
\usepackage{caption}
\usepackage{subcaption}
\usepackage{float}

\geometry{margin=2.5cm, top=3cm, bottom=3cm}

\usepackage{multirow}
% ========== Colors ==========
% Common
\definecolor{defblue}{RGB}{30, 80, 150}
\definecolor{thmgreen}{RGB}{20, 110, 60}
\definecolor{noteorange}{RGB}{200, 100, 0}
\definecolor{lightblue}{RGB}{230, 240, 255}
\definecolor{lightgreen}{RGB}{230, 255, 235}
\definecolor{lightorange}{RGB}{255, 245, 230}
\definecolor{lightgray}{RGB}{245, 245, 245}
% Extended
\definecolor{lightred}{RGB}{255, 235, 235}
\definecolor{darkred}{RGB}{180, 30, 30}
\definecolor{biogreen}{RGB}{10, 130, 80}
\definecolor{cvpurple}{RGB}{100, 40, 150}
\definecolor{injuryred}{RGB}{200, 40, 40}
\definecolor{codegray}{RGB}{240, 240, 245}
\definecolor{pipeblue}{RGB}{25, 100, 180}
\definecolor{constraintred}{RGB}{200, 50, 50}
\definecolor{termblack}{RGB}{30, 30, 30}
\definecolor{goldcolor}{RGB}{180, 140, 20}

% ========== tcolorbox environments ==========
% --- Common (all parts) ---
\newtcolorbox{keyidea}[1][]{
  colback=lightblue,
  colframe=defblue,
  fonttitle=\bfseries,
  title={Key Idea},
  #1
}

\newtcolorbox{importantnote}[1][]{
  colback=lightorange,
  colframe=noteorange,
  fonttitle=\bfseries,
  title={Important},
  #1
}

\newtcolorbox{summary}[1][]{
  colback=lightgreen,
  colframe=thmgreen,
  fonttitle=\bfseries,
  title={Summary},
  #1
}

% --- Warning/Error ---
\newtcolorbox{warningbox}[1][]{
  colback=lightred,
  colframe=darkred,
  fonttitle=\bfseries,
  title={Warning},
  #1
}

% --- Domain: Biomechanics & Clinical ---
\newtcolorbox{clinicalnote}[1][]{
  colback=lightred,
  colframe=darkred,
  fonttitle=\bfseries,
  title={Clinical Note},
  #1
}

\newtcolorbox{biomechbox}[1][]{
  colback=lightgreen,
  colframe=biogreen,
  fonttitle=\bfseries,
  title={Biomechanics},
  #1
}

\newtcolorbox{injurybox}[1][]{
  colback={injuryred!8},
  colframe=injuryred,
  fonttitle=\bfseries,
  title={Injury Risk},
  #1
}

% --- Domain: Computer Vision ---
\newtcolorbox{cvbox}[1][]{
  colback={cvpurple!5},
  colframe=cvpurple,
  fonttitle=\bfseries,
  title={Computer Vision},
  #1
}

% --- Pipeline & Setup ---
\newtcolorbox{setupbox}[1][]{
  colback=lightgray,
  colframe=gray!60,
  fonttitle=\bfseries,
  title={Setup},
  #1
}

\newtcolorbox{pipelinebox}[1][]{
  colback=pipeblue!5,
  colframe=pipeblue,
  fonttitle=\bfseries,
  title={Pipeline Step},
  #1
}

\newtcolorbox{codebox}[1][]{
  colback=codegray,
  colframe=gray!70,
  fonttitle=\bfseries\ttfamily,
  title={Code},
  #1
}

\newtcolorbox{termbox}[1][]{
  colback=termblack!5,
  colframe=termblack!60,
  fonttitle=\bfseries\ttfamily,
  title={Terminal},
  #1
}

% --- Research ---
\newtcolorbox{hypothesis}[1][]{
  colback=goldcolor!8,
  colframe=goldcolor,
  fonttitle=\bfseries,
  title={Hypothesis},
  #1
}

\newtcolorbox{resultbox}[1][]{
  colback=thmgreen!8,
  colframe=thmgreen,
  fonttitle=\bfseries,
  title={Expected Result},
  #1
}

% --- Constraints & Solutions ---
\newtcolorbox{constraintbox}[1][]{
  colback=constraintred!8,
  colframe=constraintred,
  fonttitle=\bfseries,
  title={Constraint},
  #1
}

\newtcolorbox{solutionbox}[1][]{
  colback=thmgreen!8,
  colframe=thmgreen,
  fonttitle=\bfseries,
  title={Solution},
  #1
}

\newtcolorbox{databox}[1][]{
  colback=pipeblue!5,
  colframe=pipeblue,
  fonttitle=\bfseries,
  title={Data Spec},
  #1
}

% --- Progress/Status ---
\newtcolorbox{successbox}[1][]{
  colback=lightgreen,
  colframe=thmgreen,
  fonttitle=\bfseries,
  title={Success},
  #1
}

\newtcolorbox{nextbox}[1][]{
  colback=lightorange,
  colframe=noteorange,
  fonttitle=\bfseries,
  title={Next Step},
  #1
}

% --- Fitting guide ---
\newtcolorbox{failbox}[1][]{
  colback=lightred,
  colframe=darkred,
  fonttitle=\bfseries,
  title={Failure Pattern},
  #1
}

\newtcolorbox{fixbox}[1][]{
  colback=lightgreen,
  colframe=thmgreen,
  fonttitle=\bfseries,
  title={Fix},
  #1
}

\newtcolorbox{lessonbox}[1][]{
  colback=lightorange,
  colframe=noteorange,
  fonttitle=\bfseries,
  title={Lesson Learned},
  #1
}

% --- Specification ---
\newtcolorbox{specbox}[1][]{
  colback=cvpurple!5,
  colframe=cvpurple,
  fonttitle=\bfseries,
  title={Specification},
  #1
}

\newtcolorbox{checklistbox}[1][]{
  colback=lightgreen,
  colframe=biogreen,
  fonttitle=\bfseries,
  title={Checklist},
  #1
}

% ========== Custom math commands ==========
% Vectors (lowercase)
\newcommand{\vx}{\mathbf{x}}
\newcommand{\vd}{\mathbf{d}}
\newcommand{\vo}{\mathbf{o}}
\newcommand{\vc}{\mathbf{c}}
\newcommand{\vr}{\mathbf{r}}
\newcommand{\vp}{\mathbf{p}}
\newcommand{\vv}{\mathbf{v}}
\newcommand{\vn}{\mathbf{n}}
\newcommand{\vu}{\mathbf{u}}
\newcommand{\vt}{\mathbf{t}}
\newcommand{\ve}{\mathbf{e}}
\newcommand{\vq}{\mathbf{q}}
% Vectors (uppercase)
\newcommand{\vF}{\mathbf{F}}
\newcommand{\vM}{\mathbf{M}}
\newcommand{\va}{\mathbf{a}}
\newcommand{\vg}{\mathbf{g}}
\newcommand{\vX}{\mathbf{X}}
% Bold Greek
\newcommand{\vmu}{\bm{\mu}}
\newcommand{\vbeta}{\bm{\beta}}
\newcommand{\vtheta}{\bm{\theta}}
\newcommand{\vpsi}{\bm{\psi}}
% Matrices
\newcommand{\mSigma}{\bm{\Sigma}}
\newcommand{\mR}{\mathbf{R}}
\newcommand{\mS}{\mathbf{S}}
\newcommand{\mW}{\mathbf{W}}
\newcommand{\mJ}{\mathbf{J}}
\newcommand{\mG}{\mathbf{G}}
\newcommand{\mT}{\mathbf{T}}
\newcommand{\mI}{\mathbf{I}}
\newcommand{\mB}{\mathbf{B}}
\newcommand{\mK}{\mathbf{K}}
\newcommand{\mP}{\mathbf{P}}
\newcommand{\mF}{\mathbf{F}}
\newcommand{\mE}{\mathbf{E}}
\newcommand{\mA}{\mathbf{A}}
\newcommand{\mH}{\mathbf{H}}
% Sets and groups
\newcommand{\RR}{\mathbb{R}}
\newcommand{\SO}{\mathrm{SO}}
% Units
\newcommand{\degps}{\,\text{deg/s}}
\newcommand{\Nm}{\ensuremath{\,\text{N}\cdot\text{m}}}
\newcommand{\BW}{\,\text{BW}}

% ========== Listings ==========
\usepackage{listings}

\lstset{
  basicstyle=\ttfamily\small,
  keywordstyle=\color{defblue}\bfseries,
  commentstyle=\color{thmgreen}\itshape,
  stringstyle=\color{noteorange},
  backgroundcolor=\color{codegray},
  frame=single,
  rulecolor=\color{gray!40},
  breaklines=true,
  showstringspaces=false,
  numbers=left,
  numberstyle=\tiny\color{gray},
  tabsize=4,
  captionpos=b
}

% ========== Header/Footer ==========
\pagestyle{fancy}
\fancyhf{}
\fancyhead[L]{\leftmark}
\fancyhead[R]{\thepage}
\renewcommand{\headrulewidth}{0.4pt}

% ========== Title formatting ==========
\titleformat{\chapter}[display]
  {\normalfont\huge\bfseries\color{defblue}}
  {\chaptertitlename\ \thechapter}{20pt}{\Huge}

\titleformat{\section}
  {\normalfont\Large\bfseries\color{defblue!80!black}}
  {\thesection}{1em}{}

\titleformat{\subsection}
  {\normalfont\large\bfseries\color{defblue!60!black}}
  {\thesubsection}{1em}{}


\begin{document}

% ==================== TITLE PAGE ====================
\begin{titlepage}
\centering
\vspace*{1.5cm}
{\Large\color{gray} GART-Pitch Development Log}
\vspace{0.5cm}

{\Huge\bfseries\color{defblue} 3DGS 기반 투구 동작 분석\\[0.3cm]
개발 진행 보고서}

\vspace{0.8cm}

{\LARGE\color{defblue!70!black} 카메라 보정, 세그멘테이션, 자세 추정까지\\
현재까지의 진행 상황과 다음 단계}

\vspace{2cm}

\begin{tikzpicture}[
  done/.style={rectangle, rounded corners=4pt, draw=thmgreen, fill=thmgreen!15,
    minimum width=2.5cm, minimum height=0.7cm, font=\small\bfseries, text=thmgreen!80!black, align=center},
  curr/.style={rectangle, rounded corners=4pt, draw=noteorange, fill=noteorange!15,
    minimum width=2.5cm, minimum height=0.7cm, font=\small\bfseries, text=noteorange!80!black, align=center},
  todo/.style={rectangle, rounded corners=4pt, draw=gray!50, fill=gray!10,
    minimum width=2.5cm, minimum height=0.7cm, font=\small\bfseries, text=gray, align=center},
  arrow/.style={-Stealth, thick}
]
\node[done] (cam) at (0,0) {Camera\\Calibration};
\node[done] (seg) at (3.5,0) {Segmentation\\(SAM3)};
\node[done] (op2d) at (7,0) {OpenPose\\2D \& 3D};
\node[curr] (smpl) at (10.5,0) {SMPL\\Fitting};
\node[todo] (gart) at (14,0) {GART\\Learning};

\draw[arrow, thmgreen] (cam) -- (seg);
\draw[arrow, thmgreen] (seg) -- (op2d);
\draw[arrow, noteorange] (op2d) -- (smpl);
\draw[arrow, gray!50] (smpl) -- (gart);

\node[font=\tiny\color{thmgreen}] at (0, -0.6) {완료};
\node[font=\tiny\color{thmgreen}] at (3.5, -0.6) {완료};
\node[font=\tiny\color{thmgreen}] at (7, -0.6) {완료};
\node[font=\tiny\color{noteorange}] at (10.5, -0.6) {진행 중};
\node[font=\tiny\color{gray}] at (14, -0.6) {대기};
\end{tikzpicture}

\vspace{2cm}
{\large 2026년 4월 11일}
\vspace{0.5cm}
{\large Subject 1 / Set 002}

\vspace{\fill}
{\small\color{gray} Part 1--4 교재 시리즈의 후속 실전 개발 보고서.\\
실제 데이터에 대한 카메라 보정, 세그멘테이션, 자세 추정 결과를 정리합니다.}
\end{titlepage}

\tableofcontents
\newpage


% ====================================================================
\chapter{프로젝트 개요}
% ====================================================================

\section{목표}

7대 240\,Hz RGB 카메라로 촬영한 야구 투구 동작에 대해,
GART(Gaussian Articulated Representation Templates) 기반 3D 인체 재구성을 수행하고,
Vicon 마커 기반 결과와 비교하여 마커리스 동작 분석의 정확도를 검증한다.

\section{데이터 현황}

\begin{table}[H]
\centering
\caption{보유 데이터}
\renewcommand{\arraystretch}{1.3}
\begin{tabularx}{\textwidth}{lXl}
\toprule
\textbf{데이터} & \textbf{내용} & \textbf{위치} \\
\midrule
마커리스 영상 & 7대 240\,Hz MOV (cam1: portrait, cam2--7: landscape) & \texttt{data/markerless/} \\
마커 기반 & Vicon 240\,Hz, 42 마커 Plug-in Gait, 포스플레이트 $\times$2 & \texttt{data/markerbased/} \\
세그멘테이션 & SAM3 마스크 + 배경 제거 이미지, 7대 $\times$ 1000프레임 & \texttt{data/segments/} \\
OpenPose & 2D keypoints 7대 + 3D 삼각측량 결과 & \texttt{data/openpose/} \\
카메라 보정 & DA3 PnP refined + COLMAP 형식 변환 & \texttt{data/da3\_calibration/} \\
\bottomrule
\end{tabularx}
\end{table}


\section{파이프라인 진행 상태}

\begin{table}[H]
\centering
\caption{GART-Pitch 파이프라인 진행 현황}
\renewcommand{\arraystretch}{1.4}
\begin{tabularx}{\textwidth}{clXl}
\toprule
\textbf{\#} & \textbf{단계} & \textbf{내용} & \textbf{상태} \\
\midrule
0 & 카메라 보정 & DA3 PnP refined 채택 (COLMAP 부적합 판정) & \textcolor{thmgreen}{\textbf{완료}} \\
1 & 세그멘테이션 & SAM3 7대 전체 마스크 (cam7 tracked OpenPose 재생성) & \textcolor{thmgreen}{\textbf{완료}} \\
2 & 2D 자세 추정 & OpenPose Body25 $\times$ 7카메라 & \textcolor{thmgreen}{\textbf{완료}} \\
3 & 3D 삼각측량 & DLT 기반 7뷰 삼각측량 $\rightarrow$ 25관절 3D & \textcolor{thmgreen}{\textbf{완료}} \\
4 & SMPL 피팅 & 3D keypoints $\rightarrow$ SMPL 파라미터 & \textcolor{noteorange}{\textbf{진행 중}} \\
5 & GART 학습 & 캐노니컬 Gaussian + LBS + 자세 정제 & 대기 \\
6 & 생체역학 분석 & 관절 각도, 토크, 임상 지표 & 대기 \\
7 & Vicon 검증 & MPJPE, RMSE, ICC, Bland--Altman & 대기 \\
\bottomrule
\end{tabularx}
\end{table}


% ====================================================================
\chapter{카메라 보정}
% ====================================================================

\section{두 가지 접근법 비교}

\begin{table}[H]
\centering
\caption{COLMAP vs DA3 PnP Refined 비교}
\renewcommand{\arraystretch}{1.3}
\begin{tabularx}{\textwidth}{lXX}
\toprule
\textbf{지표} & \textbf{COLMAP Round 2} & \textbf{DA3 PnP Refined} \\
\midrule
등록률 & 336/336 (100\%) & 7/7 (100\%) \\
카메라 모델 & SIMPLE\_RADIAL & PINHOLE (fx$\approx$fy) \\
$f$ 범위 & 531--1367 px (불일치) & $\sim$790 px 전체 일관 \\
카메라 배치 & GT와 불일치 (거리 오차 50\%) & 물리적 배치와 일치 \\
좌표계 & COLMAP 자체 & DA3 월드 좌표 \\
\textbf{판정} & \textcolor{darkred}{\textbf{부적합}} & \textcolor{thmgreen}{\textbf{채택}} \\
\bottomrule
\end{tabularx}
\end{table}

\begin{warningbox}
\textbf{COLMAP이 실패한 원인}:
\begin{enumerate}[nosep]
  \item \textbf{동적 장면}: 투수가 움직여서 SfM의 정적 장면 가정 위반
  \item \textbf{동일 카메라 미인식}: 7대 모두 같은 기종인데 COLMAP이 서로 다른 초점거리 추정
  \item \textbf{시야 겹침 부족}: 일부 카메라 간 특징점 매칭 실패
\end{enumerate}
\end{warningbox}

\section{DA3 PnP Refined 카메라 파라미터}

\begin{table}[H]
\centering
\caption{최종 채택된 카메라 파라미터 (DA3 PnP Refined, 원본 해상도)}
\renewcommand{\arraystretch}{1.3}
\small
\begin{tabular}{lcccccc}
\toprule
\textbf{카메라} & \textbf{해상도} & \textbf{$f_x$} & \textbf{$f_y$} & \textbf{$c_x$} & \textbf{$c_y$} & \textbf{$f_x/f_y$} \\
\midrule
cam1 & 1080$\times$1920 & 797.4 & 785.1 & 534.9 & 964.6 & 1.016 \\
cam2 & 1920$\times$1080 & 792.3 & 797.6 & 961.1 & 533.9 & 0.993 \\
cam3 & 1920$\times$1080 & 781.6 & 786.2 & 957.6 & 542.7 & 0.994 \\
cam4 & 1920$\times$1080 & 784.1 & 801.4 & 952.3 & 542.9 & 0.978 \\
cam5 & 1920$\times$1080 & 787.0 & 791.1 & 962.2 & 542.6 & 0.995 \\
cam6 & 1920$\times$1080 & 783.5 & 788.8 & 950.1 & 544.0 & 0.993 \\
cam7 & 1920$\times$1080 & 787.5 & 792.8 & 955.7 & 548.4 & 0.993 \\
\bottomrule
\end{tabular}
\end{table}

\begin{keyidea}
DA3 PnP Refined는 280$\times$280 이미지에서 보정되었다.
원본 해상도로의 역변환 과정을 정확히 추적하였다:
\begin{enumerate}[nosep]
  \item \texttt{upper\_bound\_resize}: longest side $\rightarrow$ 504 (scale=0.2625)
  \item \texttt{make\_divisible\_by\_14}: 504$\times$284 $\rightarrow$ 504$\times$280
  \item \texttt{center\_crop}: 504$\times$280 $\rightarrow$ 280$\times$280 (landscape: $c_x$ -= 112)
\end{enumerate}
이 역변환을 적용하여 COLMAP 형식(\texttt{cameras.txt}, \texttt{images.txt})으로 변환 완료.
\end{keyidea}


\section{3D 점군 (DA3 Depth Fusion)}

Depth Anything V3로 7대 카메라에서 추정한 깊이맵을 역투영하여
329,280개의 컬러 3D 점군을 생성하였다.

\begin{table}[H]
\centering
\caption{DA3 3D 점군 사양}
\renewcommand{\arraystretch}{1.3}
\begin{tabular}{ll}
\toprule
점 수 & 329,280 \\
색상 & RGB (원본 이미지에서) \\
좌표 범위 & X: $[-0.72, 1.16]$, Y: $[-0.13, 0.95]$, Z: $[-1.13, 0.96]$ \\
저장 형식 & PLY (binary, 4.7\,MB) \\
카메라 와이어프레임 & GLB에 포함 (7대) \\
\bottomrule
\end{tabular}
\end{table}


% ====================================================================
\chapter{세그멘테이션 (SAM3)}
% ====================================================================

\section{파이프라인}

\begin{cvbox}[title={SAM3 세그멘테이션 파이프라인}]
\begin{enumerate}[nosep]
  \item OpenPose Body25로 각 프레임에서 사람 검출
  \item \texttt{pitcher\_indices.json}으로 투수의 person index 식별
  \item SAM3 \texttt{set\_text\_prompt("person")} $\rightarrow$ 복수 마스크 생성
  \item 투수 키포인트와 가장 많이 겹치는 마스크 선택
  \item 바이너리 마스크 (PNG, 0/255) 저장
\end{enumerate}
\end{cvbox}


\section{cam7 문제와 해결}

\begin{warningbox}
\textbf{cam7 문제점} (기존):
\begin{itemize}[nosep]
  \item cam7 시야에 \textbf{3명}이 존재 (투수 + 앉은 사람 + 서있는 사람)
  \item OpenPose가 프레임마다 person 순서를 바꿔 \textbf{투수가 아닌 사람}이 선택됨
  \item 마스크 해상도가 \textbf{portrait (1080$\times$1920)}로 생성되어 이미지(landscape 1920$\times$1080)와 \textbf{방향 불일치}
\end{itemize}
\end{warningbox}

\begin{successbox}
\textbf{해결}:
\begin{enumerate}[nosep]
  \item \textbf{Tracked OpenPose} 적용: 프레임 간 person tracking으로 P0 = 항상 투수 보장
  \item \textbf{SAM3 재실행} (GPU 서버, A100 80GB): tracked OpenPose 기반으로 1000프레임 전체 재생성
  \item \textbf{결과}: 1000/1000 성공, 0 실패 (237초, 4.2\,fps)
  \item \textbf{해상도 정상}: 1920$\times$1080 (landscape, 이미지와 일치)
\end{enumerate}
\end{successbox}

\section{최종 세그멘테이션 결과}

\begin{table}[H]
\centering
\caption{세그멘테이션 결과 요약}
\renewcommand{\arraystretch}{1.3}
\begin{tabular}{lcccl}
\toprule
\textbf{카메라} & \textbf{마스크} & \textbf{Masked Image} & \textbf{해상도} & \textbf{상태} \\
\midrule
cam1 & 999 & 999 & 1080$\times$1920 & \textcolor{thmgreen}{정상} \\
cam2 & 999 & 999 & 1920$\times$1080 & \textcolor{thmgreen}{정상} \\
cam3 & 1000 & 1000 & 1920$\times$1080 & \textcolor{thmgreen}{정상} \\
cam4 & 1000 & 1000 & 1920$\times$1080 & \textcolor{thmgreen}{정상} \\
cam5 & 1000 & 1000 & 1920$\times$1080 & \textcolor{thmgreen}{정상} \\
cam6 & 1000 & 1000 & 1920$\times$1080 & \textcolor{thmgreen}{정상} \\
cam7 & 1000 & 1000 & 1920$\times$1080 & \textcolor{thmgreen}{\textbf{재생성 완료}} \\
\bottomrule
\end{tabular}
\end{table}


% ====================================================================
\chapter{자세 추정 (OpenPose)}
% ====================================================================

\section{2D 자세 추정}

7대 카메라에서 OpenPose Body25 (25 관절)를 사용하여 프레임별 2D 키포인트를 추출하였다.

\begin{table}[H]
\centering
\caption{2D OpenPose 결과}
\renewcommand{\arraystretch}{1.3}
\begin{tabular}{lcl}
\toprule
\textbf{카메라} & \textbf{프레임 수} & \textbf{비고} \\
\midrule
device1--6 & 999--1000 & 원본 OpenPose \\
device7 & 1000 & \textbf{Tracked} (P0=투수 보장) \\
\bottomrule
\end{tabular}
\end{table}


\section{3D 삼각측량}

\begin{cvbox}[title={DLT 삼각측량}]
7대 카메라의 2D 키포인트를 DA3 PnP Refined 카메라 파라미터로
DLT (Direct Linear Transform) 삼각측량하여 3D 관절 위치를 산출하였다.
\begin{itemize}[nosep]
  \item \textbf{알고리즘}: DLT (SVD), $P = K[R|t]$
  \item \textbf{최소 뷰}: 3대 이상에서 검출된 관절만 삼각측량
  \item \textbf{최대 재투영 오차}: 80\,px 이상은 제거
  \item \textbf{결과}: 1000 프레임 $\times$ 25 관절 $\times$ XYZ
\end{itemize}
\end{cvbox}

\begin{table}[H]
\centering
\caption{3D OpenPose 결과 --- Body25 관절 목록}
\renewcommand{\arraystretch}{1.1}
\small
\begin{tabular}{cl|cl}
\toprule
\textbf{\#} & \textbf{관절명} & \textbf{\#} & \textbf{관절명} \\
\midrule
0 & Nose & 13 & LKne \\
1 & Neck & 14 & LAnk \\
2 & RSho & 15 & REye \\
3 & RElb & 16 & LEye \\
4 & RWri & 17 & REar \\
5 & LSho & 18 & LEar \\
6 & LElb & 19 & LBToe \\
7 & LWri & 20 & LSToe \\
8 & MHip & 21 & LHeel \\
9 & RHip & 22 & RBToe \\
10 & RKne & 23 & RSToe \\
11 & RAnk & 24 & RHeel \\
12 & LHip & & \\
\bottomrule
\end{tabular}
\end{table}

\begin{successbox}
OpenPose 3D 결과를 3D 뷰어로 시각적 검증한 결과,
스켈레톤 동작이 \textbf{투구 동작과 자연스럽게 일치}하는 것을 확인하였다.
이 3D keypoints를 SMPL 피팅의 입력으로 사용한다.
\end{successbox}


% ====================================================================
\chapter{SMPL 피팅 (현재 진행 중)}
% ====================================================================

\section{기존 EasyMocap SMPL 결과 분석}

CJH4에서 수행된 EasyMocap SMPL 피팅 결과(998프레임)를 분석한 결과,
\textbf{좌표계 불일치}로 인해 스켈레톤이 올바르지 않았다.

\begin{warningbox}
\textbf{EasyMocap SMPL 실패 원인}:

\begin{table}[H]
\centering
\renewcommand{\arraystretch}{1.3}
\begin{tabular}{lcc}
\toprule
& \textbf{OpenPose 3D (MidHip)} & \textbf{SMPL (Th)} \\
\midrule
Frame 0 좌표 & $(0.12,\; 2.81,\; 1.21)$ & $(0.16,\; 0.79,\; 2.71)$ \\
\bottomrule
\end{tabular}
\end{table}

\begin{itemize}[nosep]
  \item \textbf{Y축과 Z축이 뒤바뀜}: OpenPose Y=2.81(높이) vs SMPL Y=0.79
  \item \textbf{좌표 차이 2.5}: 같은 점(골반)인데 위치가 크게 다름
  \item \textbf{글로벌 회전 과도}: Rh 평균 123$^\circ$ --- 좌표계 보정을 Rh로 흡수
  \item \textbf{원인}: EasyMocap이 2D에서 재구성하면서 다른 좌표 규약을 사용
\end{itemize}
\end{warningbox}

\section{올바른 SMPL 피팅 전략}

\begin{keyidea}
이미 정확한 \textbf{OpenPose 3D} (25관절, 1000프레임)가 있으므로,
EasyMocap(2D $\rightarrow$ 3D 재구성)을 거칠 필요 없이
\textbf{3D keypoints에서 직접 SMPL을 피팅}하는 것이 올바른 접근이다.

\begin{equation}
\hat{\vartheta}, \hat{\beta}, \hat{t} = \arg\min_{\vartheta, \beta, t}
\sum_{j \in \mathcal{M}} w_j \left\| J_j^{\text{SMPL}}(\vartheta, \beta) + t - \mathbf{p}_j^{\text{OP3D}} \right\|^2
+ \lambda_\beta \|\beta\|^2 + \lambda_\theta \|\vartheta\|^2
\end{equation}

여기서 $\mathcal{M}$은 OpenPose Body25 $\rightarrow$ SMPL 24 관절 매핑,
$w_j$는 관절별 가중치 (신뢰도 기반)이다.
\end{keyidea}

\begin{table}[H]
\centering
\caption{OpenPose Body25 $\rightarrow$ SMPL 24 관절 매핑}
\renewcommand{\arraystretch}{1.2}
\small
\begin{tabular}{cl|cl}
\toprule
\textbf{OP25} & \textbf{관절명} & \textbf{SMPL} & \textbf{SMPL 관절명} \\
\midrule
1 & Neck & 12 & Neck \\
8 & MidHip & 0 & Pelvis \\
2 & RSho & 17 & R\_Shoulder \\
3 & RElb & 19 & R\_Elbow \\
4 & RWri & 21 & R\_Wrist \\
5 & LSho & 16 & L\_Shoulder \\
6 & LElb & 18 & L\_Elbow \\
7 & LWri & 20 & L\_Wrist \\
9 & RHip & 2 & R\_Hip \\
10 & RKne & 5 & R\_Knee \\
11 & RAnk & 8 & R\_Ankle \\
12 & LHip & 1 & L\_Hip \\
13 & LKne & 4 & L\_Knee \\
14 & LAnk & 7 & L\_Ankle \\
\bottomrule
\end{tabular}
\end{table}


% ====================================================================
\chapter{프로젝트 디렉토리 구조}
% ====================================================================

\begin{lstlisting}
3dgs_to_gart_textbook/
|-- data/
|   |-- markerless/subject_1/set_002/    # 원본 MOV (7대)
|   |-- markerbased/subject_1/set_002/   # Vicon C3D + CSV
|   |-- segments/
|   |   |-- masked_images/cam1~7/        # 배경 제거 이미지
|   |   |-- masks/cam1~7/                # SAM3 바이너리 마스크
|   |-- openpose/
|   |   |-- 2d/device1~7/                # 2D keypoints JSON
|   |   |-- 3d/openpose3d_data.json      # 3D 삼각측량 결과
|   |-- da3_calibration/
|   |   |-- cameras_da3_pnp_refined.json # 원본 보정 파일
|   |   |-- camera_params.npz            # DA3 depth + intrinsics
|   |   |-- da3_pointcloud.ply           # 329K 컬러 점군
|   |   |-- colmap_format/               # 3DGS용 변환
|   |-- secret_key/                      # GPU 서버 접속 키
|
|-- colmap/
|   |-- colmap_calibration_report.pdf    # COLMAP 보정 보고서
|   |-- round2/sparse/1/                 # COLMAP 결과 (참고용)
|   |-- visualizations/                  # 3D 뷰어 HTML들
|
|-- part5_research_proposal.tex/.pdf     # 연구 제안서
|-- part5_dev_setup.tex/.pdf             # 개발 환경 가이드
|-- part5_constraints_setup.tex/.pdf     # 제한사항 분석
\end{lstlisting}


% ====================================================================
\chapter{다음 단계}
% ====================================================================

\begin{nextbox}
\textbf{즉시 필요한 작업}:

\begin{enumerate}[leftmargin=*]
  \item \textbf{SMPL 직접 피팅} (GPU 서버)\\
    OpenPose 3D (25관절) $\rightarrow$ SMPL (24관절) 최적화.
    smplx 라이브러리 + PyTorch Adam optimizer.
    프레임별 $\vartheta$ (72), $\beta$ (10), translation (3) 추정.

  \item \textbf{Vicon--RGB 시간 동기화}\\
    동작 신호 크로스-상관으로 시간 오프셋 추정.
    Vicon 240\,Hz = RGB 240\,Hz이므로 1:1 프레임 대응 가능.

  \item \textbf{GART 학습}\\
    SMPL 자세 + 7대 영상 $\rightarrow$ 캐노니컬 Gaussian + LBS 변환.
    렌더링 기반 자세 정제 (photometric loss).
\end{enumerate}
\end{nextbox}

\begin{nextbox}
\textbf{이후 단계}:
\begin{enumerate}[leftmargin=*, start=4]
  \item 생체역학 지표 산출 (ISB 관절각, MER, 외반 토크)
  \item Vicon 대비 정량적 검증 (MPJPE, RMSE, ICC, Bland--Altman)
  \item 결과 논문 작성 (Part 5 연구 제안서 기반)
\end{enumerate}
\end{nextbox}


\end{document}
